\scanpage{125}
% line head | text | 74
\mspage{74}
% line 1–5 | text | Mais l'inclusion d'un objet initial
Mais l'inclusion d'un objet initial est toujours
asphérique, et le reste par produit avec un
objet quelconque (ici $\mathord{}_{u'}\backslash^{U}$) de $\mathrm{Cat}$
(et $=$ par \unsure{tt} changement de base lisse\ldots),
d'où la conclusion. \unsure{$\backslash$}

% line 6–9 | text | Dans la suite, nous nous intéressons
Dans la suite, nous nous intéressons
\unsure{surtout|seulement} à la situation de somme amalgamée
(où on change de notations un peu, à
présent)
% line 10–12 | diagram | (*) two rows of the pushout with dotted closed sides
\diagram{\begin{tikzcd}[ampersand replacement=\&,column sep=small]
 \& U\arrow[r,hook,"i"]\arrow[d,"p"'] \& X\arrow[d,"q"'] \& T\arrow[l,dashed,hook',"j"']\arrow[d,dashed]\\
 {(*)} \& U'\arrow[r,hook,"i'"] \& X' \& {T'\simeq T}\arrow[l,dashed,hook']
\end{tikzcd}}
% line 13–22 | text | Le fermé T prime complémentaire
Le \deleted{catégorie}\inserted{fermé $T'$} complémentaire de $U'$ \inserted{dans $X'$} est alors
can. isom. au fermé $T$ compl. de $U$ dans
$X$, l'isom. étant induit par $q$. Souvent
nous identifierons $T'$ et $T$, \inserted{ce qui est licite,} en regardant
$X'$ comme déduit du recollement de
\deleted{$\operatorname{Ob} X'=$}
$U'$ et de $\unsure{t|T}$, par le bifoncteur $\mathord{U'}^{\mathrm o}\times\unsure{\bar T|T}\unsure{-|\to}\mathrm{Ens}$
explicité dans la formule $(**)$ p. 72. Nous
nous \unsure{intéressons|intéresserons} surtout à dégager des conditions
qui assurent les deux implications suivantes
% line 23–24 | display | a) and b), membership in W
\[
\begin{aligned}
\text{a)}\quad i\in W&\Longrightarrow\unsure{,}\ i'\in W\quad(\text{et }=\ \Longleftrightarrow)\\
\text{b)}\quad p\in W&\Longrightarrow q\in W\quad(\text{et }=\ \Longleftrightarrow).
\end{aligned}
\]
% archive sheet number: 125
